\section{Conclusion} \label{sec:conclusion}

Inspired by crown decompositions for \textsc{Vertex Cover} and antler decompositions for \textsc{Feedback Vertex Set}, we introduced the notion of (tight) odd cycle cuts to capture local regions of a graph in which a simple certificate exists for the membership of certain vertices in an optimal solution to \textsc{Odd Cycle Transversal}. In addition, we developed a fixed-parameter tractable algorithm to find a non-empty subset of vertices that belong to an optimal odd cycle transversal in input graphs admitting a tight odd cycle cut; the parameter~$k$ we employed is the \emph{width} of the tight \occ. Finding tight odd cycle cuts and removing the vertices certified to be in an 	optimal solution leads to search-space reduction for the natural parameterization of \textsc{Odd Cycle Transversal}. To obtain our results, one of the main technical ideas was to replace the use of minimum two-way separators that arise naturally when solving \textsc{Odd Cycle Transversal}, by minimum three-way separators that simultaneously handle breaking the odd cycles in a subgraph \emph{and} separating the resulting local bipartite subgraph from the remainder of the graph.

\subparagraph*{Theoretical challenges}
There are several interesting directions for follow-up work. We first discuss the theoretical challenges. The algorithm we presented runs in time~$2^{k^{\Oh(1)}} n^{\Oh(z)}$, where~$z$ is the order of the tight odd cycle cut in the output guarantee of \cref{thm:main}. The polynomial term in the exponent has a large degree, which is related to the size of the cut covering sets used to shrink the bipartite part of an odd cycle cut in terms of its width. While we expect that some improvements can be made by a more refined analysis, it would be more interesting to see whether an algorithmic approach that avoids color coding can lead to significantly faster algorithms. 

An odd cycle cut~$(X_B, X_C, X_R)$ of width~$|X_C| = k$ in a graph~$G$ gives rise to a $k$-secluded bipartite subgraph~$G[X_B]$; recall that a subgraph is called $k$-secluded if its open neighborhood has size~$k$. For enumerating inclusion-maximal \emph{connected} $k$-secluded subgraphs that satisfy a property~$\Pi$, a bounded-depth branching strategy was recently proposed~\cite{JansenKW23} that generalizes the enumeration of important separators. Can such branching techniques be used to improve the running time for the search-space reduction problem considered in this paper to~$2^{\Oh(k)} n^{\Oh(z)}$?

The dependence on the complexity~$z$ of the certificate is another topic for further investigation. The search-space reduction algorithm for \textsc{Feedback Vertex Set} by Donkers and Jansen~\cite{DonkersJ24} that inspired this work, also incurs a factor~$n^{\Oh(z)}$ in its running time. For \textsc{Feedback Vertex Set}, it is conjectured but not proven that such a dependence on~$z$ is unavoidable. The situation is the same for \textsc{Odd Cycle Transversal}. Is there a way to rule out the existence of an algorithm for the task of \cref{thm:main} that runs in time~$f(k, z) \cdot n^{\Oh(1)}$?

A last theoretical challenge concerns the definition of the substructures that are used to certify membership in an optimal odd cycle transversal. Our definition of an odd cycle cut~$(X_B, X_C, X_R)$ prohibits the existence of any edges between~$X_B$ and~$X_R$. Together with the requirement that~$G[X_B]$ is bipartite, this ensures that all odd cycles intersecting~$X_B$ are intersected by~$X_C$. In principle, one could also obtain the latter conclusion from a slightly less restricted graph decomposition. Since any odd cycle enters a bipartite subgraph on one edge and leaves via another, knowing that each connected component~$H$ of~$G[X_B]$ is connected to~$X_R$ by at most one edge is sufficient to guarantee that all odd cycles visiting~$X_B$ are intersected by~$X_C$. The prior work on antler structures for \textsc{Feedback Vertex Set} allows the existence of one pendant edge per component, and manages to detect such antler structures efficiently. It would be interesting to see whether our approach can be generalized for \emph{relaxed} odd cycle cuts in which each component of~$G[X_B]$ has at most one edge to~$X_R$. To adapt to this setting, one would have to refine the type of three-way separation problem that is used in the graph reduction step. 

For \textsc{Odd Cycle Transversal}, one could relax the definition of the graph decomposition even further: to ensure that odd cycles visiting~$X_B$ are intersected by~$X_C$, it would suffice for each connected component~$H$ of~$G[X_B]$ to have at most one neighbor~$v_H$ in~$X_R$, as long as all vertices of~$H$ adjacent to~$v_H$ belong to the \emph{same} side of a bipartition of~$H$.

\subparagraph*{Practical challenges}
Since the investigation of search-space reduction is inspired by practical considerations, we should not neglect to discuss practical aspects of this research direction. While we do not expect the algorithm as presented here to be practical, it serves as a proof of concept that rigorous guarantees on efficient search-space reduction can be formulated. Our work also helps to identify the types of substructures that can be used to reason locally about membership in an optimal solution. Apart from finding faster algorithms in theory and experimenting with their results, one could also target the development of specialized algorithms for concrete values of~$k$ and~$z$. 

For~$k=1$, a tight odd cycle cut of width~$1$ effectively consists of a cutvertex~$c$ of the graph whose removal splits off a bipartite connected component~$B$ but for which the subgraph induced by~$B \cup \{v\}$ contains an odd cycle. Preliminary investigations suggest that in this case, an algorithm that computes the block-cut tree, analyzes which blocks form non-bipartite subgraphs, and which cut vertices break all the odd cycles in their blocks, can be engineered to run in time~$\Oh(|V(G)| + |E(G)|)$ to find a vertex~$v$ belonging to an optimal odd cycle transversal when given a graph that has a tight odd cycle cut of width~$k=1$. Do linear-time algorithms exist for~$k>1$? These would form valuable reduction steps in algorithms solving \textsc{Odd Cycle Transversal} exactly, such as the one developed by Wernicke~\cite{Wernicke14}.

The~$k=1$ case of the \emph{relaxed} odd cycle cuts described above are in fact used as one of the reduction rules in Wernicke's algorithm~\cite[Rule 7]{Wernicke14}. His reduction applies whenever there is a triangle~$\{u,v,w\}$ in which~$w$ has degree two and~$v$ has degree at most three. Under these circumstances, there is an optimal solution that contains~$u$ while avoiding~$v$ and~$w$: since the removal of~$u$ decreases the degree of~$w$ to one, while~$w$ is one of the at most two remaining neighbors of~$v$, the removal of~$u$ breaks all odd cycles intersecting~$\{u,v,w\}$. This corresponds to the fact that the triple~$(X_B = \{v,w\}, X_C = \{u\}, X_R = V(G) \setminus \{u,v,w\})$ forms a tight \emph{relaxed} odd cycle cut. We interpret the fact that the~$k=1$ case was developed naturally in an existing algorithm as encouraging evidence that refined research into search-space reduction steps can eventually lead to impact in practice.